\documentclass[10pt,a4paper]{amsart}
\usepackage[margin=19mm]{geometry}
\usepackage{amsmath,amssymb,mathtools}
\usepackage[scr=rsfs,cal=euler]{mathalfa}
\usepackage{xfrac}
\usepackage[unicode,hidelinks,bookmarks=false]{hyperref}
\newcommand{\ie}{i.e.\ }
\newcommand{\cf}{cf.\ }
\newcommand{\resp}{respectively\ }
\newcommand{\Subsection}{\subsection}
\providecommand{\nobreakdash}{\nobreak\hbox{-}}
\setlength{\emergencystretch}{2em}
\multlinegap0pt
\usepackage{bm}
\usepackage{bbm}
\usepackage{dsfont}
\usepackage{xspace}
\usepackage{cancel}
\usepackage{mathtools}
\usepackage{cleveref}
\usepackage{amsthm}
\makeatletter
\@ifundefined{lemma}{\newtheorem{lemma}{Lemma}}{}
\@ifundefined{proposition}{\newtheorem{proposition}{Proposition}}{}
\makeatother
\makeatletter
\expandafter\ifx\csname customlemma\endcsname\relax
\newenvironment{customlemma}[1]
  {\renewcommand\theinnercustomlemma{#1}\innercustomlemma}
  {\endinnercustomlemma}
\fi
\expandafter\ifx\csname customthm\endcsname\relax
\newenvironment{customthm}[1]
  {\renewcommand\theinnercustomthm{#1}\innercustomthm}
  {\endinnercustomthm}
\fi
\makeatother
\title[PAC-Bayes Bounds on Variational Tempered Posteriors for Mark]{PAC-Bayes Bounds on Variational Tempered Posteriors for Markov Models}
\author{Imon Banerjee, Vinayak A. Rao, Harsha Honnappa}
\date{Source: arXiv:2101.05197v1 (CC BY 4.0)}
\begin{document}
\maketitle

\noindent\emph{Source and licence.} PAC-Bayes Bounds on Variational Tempered Posteriors for Markov Models. Version arXiv:2101.05197v1 (CC BY 4.0), distributed under the Creative Commons Attribution 4.0 International licence, as recorded in the arXiv licence metadata for this exact version and shown on the abstract page https://arxiv.org/abs/2101.05197v1. Full rights notice: licenses/arxiv-2101.05197-CC-BY-4.0.txt.\par\smallskip

\noindent\textit{Editorial scope.} Counted unit: the Proposition whose source label is \texttt{lem:kl-prior} in the article (source0.tex lines 779--781, source label \texttt{lem:kl-prior}), delivered together with the article's own complete original proof of it (source0.tex lines 783--837, one \texttt{proof} environment of 55 lines). No mathematics is written, repaired or re-derived: statement and proof are the article's own text line for line. Required context retained uncounted: prop:kl-prior-unif (lines 719--721). Every definition, display and label of those lines is retained unchanged. Omitted from this package and not redistributed: 1--718, 722--778, 782--782, 838--1913. Nothing inside a retained excerpt is omitted or altered: every retained body is delivered exactly as the article's lines are written, so no omission receipt of any kind accompanies this package. Citations: the retained excerpts carry no citation command at all, so no bibliography entry of the article is transcribed and no reference is redistributed. Reference adaptation: none was needed and none was made; every reference inside the delivered unit resolves inside it, so no unresolved reference or citation is produced. Printed slips kept literal: no printed word, symbol, hypothesis or number of the counted statement or of its proof was altered. Layout and class: the article's own class is \texttt{amsart} with packages including inputenc, ulem, authblk, babel, amssymb, geometry, verbatim, amsthm, amsmath, amsaddr, graphicx, todonotes, hyperref, comment; the wrapper is the lane's pinned \texttt{amsart} editorial wrapper at 10pt on A4, which restates the theorem-like environments the retained text needs and adds the article's own macro definitions that the retained text needs (none), with extra wrapper packages (bm, bbm, dsfont, xspace, cancel, mathtools, cleveref). The delivered numbering is the unit's own. Assets: no retained span contains a figure environment, a graphics-inclusion command, a PDF-page inclusion, a table or a TikZ picture; no image file is copied into this package, the package declares no asset dependency and no image file is redistributed with it. Licence: version arXiv:2101.05197v1 of this article is distributed under the Creative Commons Attribution 4.0 International licence, as recorded in the arXiv licence metadata for this exact version and shown on the abstract page https://arxiv.org/abs/2101.05197v1. A full rights notice is delivered at licenses/arxiv-2101.05197-CC-BY-4.0.txt. Characters: every retained excerpt is pure ASCII, so the delivered engine, XeLaTeX with its default Latin Modern text font, reports no missing character.
\begin{lemma}~\label{prop:kl-prior-unif}
Let $\theta_0\in(0,1)$. Let, $\rho_n$ be a sequence of Beta distributions with parameters $\alpha_n=n\theta_0$ and $\beta_n=n(1-\theta_0)$. Let $\pi$ denote an uniform distribution, $U(0,1)$. Then, $\mathcal{K}(\rho_n,\pi)<C+\frac{1}{2}\log (n)$, for some constant $C>0$. 
\end{lemma}

\begin{proposition}~\label{lem:kl-prior}
Let $\theta_0\in(0,1)$. Let, $\rho_n$ be a sequence of Beta distributions with parameters $\alpha_n=n\theta_0$ and $\beta_n=n(1-\theta_0)$. Let $\pi$ denote an Beta distribution, with parameters $(\alpha,\beta)$. Then, $\mathcal{K}(\rho_n,\pi)<C+\frac{1}{2}\log(n)$, for some constant $C>0$. 
\end{proposition}

\begin{proof}
The KL-divergence between $\rho_n$ and $\pi$ can be written as,
\begin{align*}
    \mathcal{K}(\rho_n,\pi)= & \int \log\left(\frac{\rho_n}{\pi}\right)\rho_n(d\theta)=\int \log\left(\frac{\rho_n}{U}\right)\rho_n(d\theta)+\int \log\left(\frac{U}{\pi}\right)\rho_n(d\theta),
\end{align*}
where, $U$ is an uniform distribution on $(0,1)$.
%
We analyze the second term in the above expression. The second term can be written as,
\begin{align*}
    \int \log\left(\frac{U}{\pi}\right)\rho_n(d\theta) & = \int \log \left( \frac{1}{\frac{1}{\text{Beta}(\alpha,\beta)}\theta^{\alpha-1}(1-\theta)^{\beta-1}} \right)\rho_n(d\theta)\\
     & = C_1-(\alpha-1)\int \log(\theta) \rho_n(d\theta)-(\beta-1)\int \log(1-\theta) \rho_n(d\theta),
     \intertext{where $C_1$ is $\log(\text{Beta}(\alpha,\beta))$. Since, $\rho_n$ follows a Beta distribution with parameters $\alpha_n=n\theta_0$ and $\beta_n=n(1-\theta_0)$, we get that,}
     \int \log\left(\frac{U}{\pi}\right)\rho_n(d\theta)& =
     C_1-(\alpha-1) \left[\psi(\alpha_n)-\psi(\alpha_n+\beta_n)\right]-(\beta-1)\left[\psi(\beta_n)-\psi(\alpha_n+\beta_n)\right]
\end{align*}
Since, $\log(x)-\frac{1}{x}<\psi(x)<\log(x)-\frac{1}{2x}$, looking at the term $\left[\psi(\alpha_n)-\psi(\alpha_n+\beta_n)\right]$, we get that, 
\begin{align*}
    \left[\psi(\alpha_n)-\psi(\alpha_n+\beta_n)\right] & = \left[\psi(n\theta_0)-\psi(n\theta_0+n(1-\theta_0))\right]\\
    & = \left[\psi(n\theta_0)-\psi(n)\right].\\
    \intertext{Using the upper bound on $\psi(n\theta_0)$ and the lower bound on $\psi(n)$, we get}
    \left[\psi(n\theta_0)-\psi(n)\right] & < \log(n\theta_0)-\frac{1}{2n\theta_0}-\log(n)+\frac{1}{n\theta_0}\\
    & = \log(\theta_0)+\frac{1}{2n\theta_0}.
    \intertext{We can also get a lower bound very similarly.}
    \left[\psi(\alpha_n)-\psi(\alpha_n+\beta_n)\right] & >\log(\theta_0)-\frac{1}{2n\theta_0}.
\end{align*}
Therefore it follows that
\begin{equation}
    \log(\theta_0)-\frac{1}{2n\theta_0}<\left[\psi(\alpha_n)-\psi(\alpha_n+\beta_n)\right]< \log(\theta_0)+\frac{1}{2n\theta_0}.
\end{equation}
Similarly, we have
\begin{equation}
    \log(1-\theta_0)-\frac{1}{2n(1-\theta_0)}<\left[\psi(\beta_n)-\psi(\alpha_n+\beta_n)\right]< \log(1-\theta_0)+\frac{1}{2n(1-\theta_0)}.
\end{equation}
Consequently, for a large enough constant $C_2>0$ it follows that
\begin{align*}
    -C_2<\left[\psi(\alpha_n)-\psi(\alpha_n+\beta_n)\right]< C_2,
\end{align*}
and
\begin{align*}
    -C_2<\left[\psi(\beta_n)-\psi(\alpha_n+\beta_n)\right]< C_2.
\end{align*}
Without loss of generality assume that $\min\{\alpha-1,\beta-1\}>0$. Then we get,
\begin{align*}
    -(\alpha-1)C_2<(\alpha-1)\left[\psi(\alpha_n)-\psi(\alpha_n+\beta_n)\right]< (\alpha-1)C_2,
\end{align*}
and
\begin{align*}
    -(\alpha-1)C_2<(\alpha-1)\left[\psi(\beta_n)-\psi(\alpha_n+\beta_n)\right]< (\alpha-1)C_2.
\end{align*}
If, either of $\alpha-1$ or $\beta-1$ was negative, then the direction of inequality on the previous equation would have reversed. However, due to the symmetric nature of the inequality, that would not have made a difference. Using the above bounds we finally show that, 
\begin{align*}
    C_1-(\alpha-1) \left[\psi(\alpha_n)-\psi(\alpha_n+\beta_n)\right]-(\beta-1)\left[\psi(\beta_n)-\psi(\alpha_n+\beta_n)\right]<C_1+C_2\theta_0+C_2(1-\theta_0)<C,
\end{align*}
for some large constant C. Finally, we upper bound $\int \log\left(\frac{\rho_n}{U}\right)\rho_n(d\theta)$ by \Cref{prop:kl-prior-unif} thereby completing the proof.
\end{proof}

\end{document}
